\section{Functions}

\begin{frame}[fragile]{Returning structs from functions}
  In the following code example we have a copyable \code{IntBox} class.\\

  \begin{minted}{cpp}
    /* 3.02 */
    IntBox sum(IntBox& a, IntBox& b) {
        auto result = *(a.mData) + *(b.mData);
        return IntBox(result);
    }
    IntBox product(IntBox& a, IntBox& b) {
        auto result = *(a.mData) * *(b.mData);
        auto retval = IntBox{result};
        return retval;
    }
    int demo() {
      auto a = IntBox{5}; auto b = IntBox{6};
      auto sumResult = sum(a, b);
      auto prodResult = product(a, b);
    }
  \end{minted}
\end{frame}

\begin{frame}[fragile]{Returning structs from functions}
  In the stack frame of the caller (\code{demo()}), the compiler allocates a \emph{return\_slot}.
  The address of the return slot is implicitly passed as an argument.
  \begin{itemize}
    \item if the returned expr is a local variable, it gets copy-ctored into the return slot.
    \item if the returned expr is a ctor call, the ctor is run directly on the return slot.
  \end{itemize}

  Conceptually, this is what happens under the hood:
  \begin{minted}{cpp}
/* 3.01 */
auto a = IntBox{5}; auto b = IntBox{6};
alignas(IntBox) uint8_t _storageSumResult[sizeof(IntBox)]; // allocate return slot on stack
// inlined function
{   auto* _returnSlot = reinterpret_cast<IntBox *>(_storageSumResult);
    auto& argA = a; auto& argB = b; // argument binding
{
    auto result = *(argA.mData) + *(argB.mData);
    new (_returnSlot) IntBox(result); // RVO: constructor runs directly in return slot
}}
  \end{minted}
\end{frame}

\begin{frame}[fragile]{Returning structs from functions}
  \includesvg[inkscapelatex=false, width=\textwidth]{assets/return_slots.drawio.svg}
\end{frame}

\begin{frame}[fragile]{Temporaries}
  If the caller doesn't bind the return value to a variable,
  it still allocates a return slot on the stack.\\
  Methods may be called on this temporary while within the full expression.\\
  However immediately following the full expression, the temporary is destroyed.
  \begin{minted}{cpp}
    /* 3.01 */
    int demo() {
        // return slot is to a temporary, rather than an IntBox object
        std::cout << product(a, b).get() << std::endl;
        // _tmp.~IntBox() // the compiler generates a call to the destructor here
    }
  \end{minted}

  It is possible to bind the temporary to an r-value reference (later section),
  for example by passing it into a function that takes an r-value reference as its argument.
\end{frame}

\begin{frame}[fragile]{Issue: Unnecessary Copies}
  Our \code{IntBox} from 3.01 is copyable, but has no move constructor.
  \begin{minted}{cpp}
    /* 3.01 */
    IntBox product(IntBox& a, IntBox& b) {
        auto result = *(a.mData) * *(b.mData);
        auto retval = IntBox{result}; // allocates an int
        return retval;                // copy ctor: allocates another int
    }                                 // ~retval(): frees the first int

    int demo() {
        auto a = IntBox{5}; auto b = IntBox{6};
        auto c = product(a, b);       // c is the return slot
    }
  \end{minted}
  Without RVO, \code{retval} is copy constructed into the return slot,
  then immediately destroyed.
  We pay for an extra allocation and copy, only to throw the original away.\\[0.5em]
  \textcolor{sqt-question}{\code{retval} is about to die anyway.
  Can we steal its resources instead of copying them?}
\end{frame}

\begin{frame}[fragile]{Solution: Define a Move Constructor}
  \begin{minted}{cpp}
    IntBox(IntBox&& other) : mData{std::exchange(other.mData, nullptr)} {}
    IntBox product(IntBox& a, IntBox& b) {
        auto result = *(a.mData) * *(b.mData);
        auto retval = IntBox{result};
        return retval; // move construct into returnSlot: no allocation, very fast
    }                  // ~retval(): destroys moved-from object
  \end{minted}

  \begin{sqt:result}
    A local returned by name is treated as an rvalue.
    The copy ctor was previously called, because const lrefs can also bind to rvalues.
  \end{sqt:result}

  \begin{sqt:tip}
    Don't write \code{return std::move(retval);}, as this disables NRVO.
  \end{sqt:tip}
\end{frame}